\section*{第 \ref{ch:b2:thermal-radiation} 章——热辐射}

\begin{solution}{exo:b2:thermal-radiation:1}
$\lambda_{\max}$: \qty{0.50}{\micro m}, \qty{9.4}{\micro m}, \qty{1.04}{\micro m},
\qty{1.06}{mm}, \qty{38}{\micro m}, \qty{3.2}{\micro m}; $\sigma T^4$: \qty{6.4e7}{},
\qty{520}{}, \qty{3.5e6}{}, \qty{3.1e-6}{}, \qty{2.0}{}, \qty{3.7e4}{W/m^2}。可见：
太阳与灯丝；火钳发出暗红色的光，
其短波尾部仅占其功率的万分之一。
\end{solution}

\begin{solution}{exo:b2:thermal-radiation:2}
$S = \sigma T_\odot^4(R_\odot/d)^2 = \qty{1360}{W/m^2}$; $L = 4\pi d^2S =
\qty{3.8e26}{W}$; $\pi R^2S = \qty{1.7e17}{W}$; \qty{2e7}{W} 每人——九
千倍于世界的消耗量。
\end{solution}

\begin{solution}{exo:b2:thermal-radiation:3}
发出 \qty{830}{W}，吸收 \qty{700}{W}，净损失 \qty{130}{W}；$h_{\text{r}} =
4\varepsilon\sigma T_0^3 = \qty{5.6}{W/m^2/K}$。超过人体的代谢功率：
衣物可减少这部分损失；墙壁温度为 \qty{15}{\celsius} 时，
又增加约 \qty{50}{W} 的损失——人在暖空气中却身处冷墙之间，会感到寒冷。
\end{solution}

\begin{solution}{exo:b2:thermal-radiation:4}
$A = P/\varepsilon\sigma T^4 = \qty{49}{mm^2}$; $\ell = A/\pi d = \qty{31}{cm}$;
\qty{4.8}{W} 的光；一只 \qty{16}{W} 的发光二极管。
\end{solution}

\begin{solution}{exo:b2:thermal-radiation:5}
(a) $\eu^x - 1 \approx x$：$u_\nu = (8\pi\nu^2/c^3)\,k_BT$——每个模式 $k_BT$。
(b) $\int\nu^2\dd\nu$ 发散。(c) $\nu = k_BT/h = \qty{6e12}{Hz}$（\qty{50}{\micro m}）：
在 \qty{300}{K} 时，无线电波与微波深处于经典区（
电阻器噪声为每单位带宽 $k_BT$）。(d) 找到激发量子 $h\nu$ 的概率，
在 $T$ 下即为玻尔兹曼因子。
\end{solution}

\begin{solution}{exo:b2:thermal-radiation:6}
(a) $x = 4.965$. (b) $hc/4.965k_B = \qty{2.90e-3}{m.K}$. (c) $x^3/(\eu^x - 1)$
在 $3(1 - \eu^{-x}) = x$ 处取峰值 $2.82$；按单位频率的
与按单位波长的能量密度为不同的函数（$\dd\nu = c\,\dd\lambda/\lambda^2$），
并在不同位置达到峰值。(d) \qty{10000}{K}；偏蓝的白色。
\end{solution}

\begin{solution}{exo:b2:thermal-radiation:7}
(a) $u = (8\pi^5k_B^4/15h^3c^3)T^4$, $\sigma = cu/4T^4 = \qty{5.67e-8}{W/m^2/K^4}$.
(b) 各向同性与 $\cos\theta$ 投影：$\int\cos\theta\,\dd\Omega/4\pi = 1/4$
对半球面积分。(c) $n = 8\pi \times 2.404\,(k_BT/hc)^3 = 60.4\,(k_BT/hc)^3$；
平均能量 $(\pi^4/15)/2.404 = 2.70\,k_BT$。(d) $5 \times 10^8$ 与 $400$
每立方厘米。
\end{solution}

\begin{solution}{exo:b2:thermal-radiation:8}
(a) $\sigma(T_1^4 - T_2^4)$。(b) $T_{\text{s}}^4 = (T_1^4 + T_2^4)/2$；减半。(c)
$1/(n + 1)$。(d) 对各次反射求和，$\sigma(T_1^4 - T_2^4)/(2/\varepsilon -
1)$；$\varepsilon = 0.05$：$\qty{460}{W/m^2}/39 = \qty{12}{W/m^2}$——一只表面积
为 \qty{0.05}{m^2} 的杜瓦瓶损失 \qty{0.6}{W}，每天蒸发
四分之一千克氮气。
\end{solution}

\begin{solution}{exo:b2:thermal-radiation:9}
(a) $\varepsilon\sigma(T^4 - 250^4) = h(278 - T)$: $T \approx \qty{266}{K}$. (b)
\qty{277}{K}。(c) 较大的 $h$ 使地表温度紧贴气温。(d) 林冠
温度为 \qty{278}{K}，取代了 \qty{250}{K} 的天空。
\end{solution}

\begin{solution}{exo:b2:thermal-radiation:10}
(a) \qty{231}{K}, \qty{255}{K}, \qty{210}{K}. (b) \qty{506}{K}, \qty{33}{K}, \qty{5}{K}.
(c) 每一不透明层均向上、向下辐射；由平衡关系给出
$T_k^4 = (k + 1)T_{\text{e}}^4$ 自顶层推得：对金星 $n \approx 100$。(d)
$((1 - A)S/\sigma)^{1/4} = \qty{381}{K}$；无大气时，
热量无法输送到夜侧，月壤储热也很少：地表
约一天之内即将白天所获热量辐射殆尽，并稳定在
\qty{100}{K} 附近。
\end{solution}

\begin{solution}{exo:b2:thermal-radiation:11}
(a) $4\sigma T_{\text{e}}^3 = \qty{3.8}{W/m^2/K}$. (b) \qty{1.0}{K}. (c) $C =
\qty{2.1e8}{J/m^2/K}$，$\tau = C/4\sigma T_{\text{e}}^3 \approx \qty{5.6e7}{s}$，即两
年。(d) 变暖的空气可容纳更多水蒸气（一种温室气体），并
使冰融化（反照率降低）：两者均会放大最初的增温。
\end{solution}

\begin{solution}{exo:b2:thermal-radiation:12}
(a)
\[
\frac{M_1}{M_2} = \Big(\frac{\lambda_2}{\lambda_1}\Big)^5\exp\Big[-\frac{hc}{k_BT}\Big(\frac{1}{\lambda_1} - \frac{1}{\lambda_2}\Big)\Big] :
\]
$\varepsilon$ 约去。(b) $T = \qty{14388}{\micro m.K} \times 0.427/(5\ln 1.385 +
\ln 10) = \qty{1570}{K}$。(c) 接收量 $\propto 0.1 \times 350 + 0.9 \times 293 =
299$；报告值 $(299 - 0.05 \times 293)/0.95 \approx \qty{299}{K}$：漏掉了
\qty{350}{K}。(d) 其\omterm{prop:b2:thermal-radiation:kirchhoff}{发射率}即为所假设的值。
\end{solution}

\begin{solution}{pb:b2:thermal-radiation:1}
\textbf{1.} \qty{5800}{K}.

\textbf{2.} \qty{6.4e7}{W/m^2}; \qty{3.9e26}{W}.

\textbf{3.} $L/4\pi d^2 = \qty{1.37e3}{W/m^2}$; $u = S/c = \qty{4.6e-6}{J/m^3}$;
$P = S/c = \qty{4.6e-6}{Pa}$（若被吸收）；若被反射则为两倍。

\textbf{4.} $2.7k_BT = \qty{2.2e-19}{J}$ (\qty{1.35}{eV}); $6 \times 10^{21}$
个光子，每秒每平方米。

\textbf{5.} $3.9 \times 10^{26} \times 1.4 \times 10^{17} = \qty{5.5e43}{J}$; $Mc^2 =
\qty{1.8e47}{J}$：相当于其质量的万分之三。

\textbf{6.} \qty{1.2e17}{W}; $T_{\text{e}} = \qty{255}{K}$.

\textbf{7.} 在 $\pi R^2$ 上截获，经 $4\pi R^2$ 辐射出去，借助于
自转与风；仅就日下点而言：$((1 - A)S/\sigma)^{1/4} =
\qty{360}{K}$.

\textbf{8.} \qty{381}{K}；夜间储存的热量（$\sim C\Delta T$）通过
辐射散失，时间常数 $C/4\sigma T^3 \approx \qty{1}{day}$，与两周的黑夜相比很短：
地面温度降至其辐射恰好与来自下方的微弱热流相平衡之处，
即 \qty{100}{K}。

\textbf{9.} \qty{10}{\micro m}，是太阳的二十倍。

\textbf{10.} $(0.65/0.70)^{1/4} = 0.982$：降低 \qty{4.7}{K}，为 \qty{250}{K}。

\textbf{11.} 空间：$\sigma T_{\text{a}}^4 = \sigma T_{\text{e}}^4$；大气层：$\sigma
T_{\text{s}}^4 = 2\sigma T_{\text{a}}^4$；地表：$\sigma T_{\text{s}}^4 = \sigma
T_{\text{e}}^4 + \sigma T_{\text{a}}^4$。由此可得 $T_{\text{s}} = 2^{1/4}T_{\text{e}} = \qty{303}{K}$。

\textbf{12.} $T_{\text{s}}^4 = T_{\text{e}}^4/(1 - \varepsilon/2)$；$(288/255)^4 = 1.63$
可得 $\varepsilon = 0.77$。

\textbf{13.} 水蒸气、二氧化碳、甲烷、臭氧在红外波段吸收
（振动谱带：H$_2$O 位于 \qty{6}{\micro m} 附近及
\qty{17}{\micro m} 以外，CO$_2$ 位于 \qty{15}{\micro m}）；N$_2$ 和 O$_2$ 不吸收（无
偶极矩）；窗口为 \qty{8}{}--\qty{13}{\micro m}。

\textbf{14.} $T_{\text{a}}^4 = T_{\text{s}}^4/2$：$\varepsilon\sigma T_{\text{a}}^4 \approx
\qty{150}{W/m^2}$，对比被吸收的太阳光 \qty{238}{W/m^2}（
多层大气回送的更多，约 \qty{330}{W/m^2}）。

\textbf{15.} $3.7/3.8 \approx \qty{1}{K}$。

\textbf{16.} \qty{2.1e8}{J/m^2/K}；$\tau \approx \qty{5.6e7}{s}$，约两年。

\textbf{17.} 云层反射太阳光，并向下辐射红外线；夜间
只有后者起作用：多云的夜晚更暖和。

\textbf{18.} \qty{49}{mm^2}；\qty{31}{cm}；卷成螺旋形以便装入，并
减少向填充气体的对流散失。

\textbf{19.} \qty{1.04}{\micro m}；\qty{4.8}{W}；$8\%$。

\textbf{20.} $R = V^2/P = \qty{880}{\ohm}$；冷态电阻仅 \qty{88}{\ohm}，电流为
\qty{2.6}{A}，而非 \qty{0.26}{A}——灯泡在接通瞬间即烧坏。

\textbf{21.} $49 \times (2800/3200)^4 = \qty{29}{mm^2}$；可见光 \qty{8.4}{W}；蒸发速率经八次
倍增，寿命缩短 $250$ 倍——仅为数小时。

\textbf{22.} 蒸发出的钨与卤素结合，并被送回
热灯丝上，而不是使玻壳发黑：灯丝可在
正常寿命下工作得更热，外壳为小型石英泡壳。

\textbf{23.} \qty{55}{W} 以红外辐射及经气体的传导形式散失；
玻璃吸收了一部分红外线。

\textbf{24.} $C = \qty{2.8e-3}{J/K}$；由 $P \propto T^4$，使 $T$ 减半所需时间
为 $7CT_0/3P \approx \qty{0.3}{s}$——而辉光随电网
以 \qty{0.04}{s} 的滞后变化，这也是它几乎不闪烁的原因。

\textbf{25.} 温度决定光谱的形状与峰值
（$\lambda_{\max}T$），以及总辐射（$\sigma T^4$）；灯丝、太阳与地球
是三个\omterm{def:b2:thermal-radiation:blackbody}{黑体}，温度分别为 \qty{2800}{}、\qty{5800}{} 和 \qty{288}{K}，每个都
使所吸收的能量与 $\sigma T^4$ 相平衡。
\end{solution}
